\section{Exercises}\label{sec:07-group-theorems:05-exercises:1}

\textbf{Key points.} Three facts now hold for \emph{every} group: the identity is unique, a left inverse equals the (unique) two-sided inverse, and \((ab)^{-1}=b^{-1}a^{-1}\). The recurring proof pattern is relating both sides of a goal to a common third expression, or padding with the identity and cancelling; once a uniqueness fact is proved, later goals can be \emph{characterized} by it instead of computed directly.

\begin{enumerate}
\def\labelenumi{\arabic{enumi}.}
\tightlist
\item
  Prove \texttt{theorem inv\_inv (a : G) : Grp.inv (Grp.inv a) = a}. Before writing any tactics, consider: does this match the shape of a lemma already in hand (Theorem 2 again)? What single fact about \texttt{a} and \texttt{Grp.inv a} would permit invoking it directly?
\item
  Prove \texttt{theorem cancel\_left (a b c : G) (h : Grp.op a b = Grp.op a c) : b = c}. Strategy hint: \texttt{b} and \texttt{c} cannot be rewritten directly in isolation. Instead, apply \texttt{Grp.op (Grp.inv a)} to \emph{both sides} of \texttt{h} first (as a \texttt{have}), then simplify each side using \texttt{assoc}/\texttt{inv\_left}/\texttt{id\_left}, the same ``regroup, then cancel'' pattern as Theorem 3.
\end{enumerate}

Solutions: \hyperref[sec:14-appendix-solutions:06-chapter-7:1]{Appendix, Chapter 7}.

\section{Next}\label{sec:07-group-theorems:05-exercises:2}

Continue to \hyperref[sec:08-rings:00-index:1]{Chapter 8: Rings}.
